% !TEX program = pdflatex
% FRE6883 Recitation 1 -- Opening and diagnostic
\documentclass[10pt,aspectratio=169]{beamer}

% Language & encoding
\usepackage[utf8]{inputenc}
\usepackage[T1]{fontenc}
\usepackage[english]{babel}
\usepackage{lmodern}
\usepackage{microtype}

% Math / tables / layout
\usepackage{amsmath,amssymb,bm,mathtools}
\usepackage{booktabs}
\usepackage{multicol}
\usepackage{changepage}

% Graphics
\usepackage{graphicx}
\usepackage{tikz}
\usetikzlibrary{arrows.meta,positioning,calc,decorations.pathreplacing}
\usepackage{pgfplots}
\pgfplotsset{compat=1.18}

% Theme (Metropolis)
\usetheme[numbering=fraction,progressbar=frametitle]{metropolis}
\setbeamertemplate{frame numbering}[fraction]
\setbeamertemplate{navigation symbols}{}

% Bootcamp palette
\definecolor{BootcampBlueDark}{HTML}{1A3C68}
\definecolor{BootcampBlue}{RGB}{31,110,178}
\definecolor{AccentGold}{HTML}{DAA520}
\definecolor{SoftGray}{RGB}{245,246,248}
\definecolor{TextBlack}{HTML}{000000}
\definecolor{Crimson}{HTML}{990000}
\definecolor{MatrixGreen}{HTML}{228B22}
\definecolor{MatrixRed}{HTML}{B22222}

% Global formatting
\setbeamercolor{normal text}{fg=TextBlack,bg=white}
\setbeamercolor{title}{fg=TextBlack}
\setbeamercolor{structure}{fg=BootcampBlueDark}
\setbeamercolor{progress bar}{fg=BootcampBlueDark,bg=gray!30}

\setbeamercolor{frametitle}{bg=BootcampBlueDark,fg=white}
\setbeamerfont{frametitle}{series=\bfseries,size=\large}
\setbeamertemplate{frametitle}{
  \nointerlineskip
  \begin{beamercolorbox}[wd=\paperwidth,ht=3.2ex,dp=1.4ex,leftskip=1em,rightskip=1em]{frametitle}
    \usebeamerfont{frametitle}\insertframetitle
  \end{beamercolorbox}
}

\setbeamercolor{block title example}{bg=BootcampBlueDark!90!black,fg=white}
\setbeamercolor{block body example}{bg=BootcampBlueDark!5!white,fg=black}
\setbeamercolor{block title proposition}{bg=MatrixGreen!75!black,fg=white}
\setbeamercolor{block body proposition}{bg=MatrixGreen!7!white,fg=black}
\setbeamercolor{block title illustration}{bg=AccentGold!85!black,fg=white}
\setbeamercolor{block body illustration}{bg=AccentGold!10!white,fg=black}
\setbeamercolor{block title proof}{bg=gray!70!black,fg=white}
\setbeamercolor{block body proof}{bg=SoftGray,fg=black}
\setbeamertemplate{blocks}[rounded][shadow=false]

\newenvironment{definitionblock}[1]{%
  \par\vspace{0.4em}%
  \begin{beamercolorbox}[
      wd=\textwidth,sep=0.4em,leftskip=0.3em,rightskip=0.6em,
      rounded=false,shadow=false
    ]{block title example}%
    \raisebox{0.15em}{\usebeamerfont{block title example}#1}%
  \end{beamercolorbox}%
  \vspace{-0.4em}%
  \begin{beamercolorbox}[
      wd=\textwidth,sep=0.9em,leftskip=0.6em,rightskip=0.6em,
      rounded=false,shadow=false
    ]{block body example}%
    \usebeamerfont{block body example}%
}{%
  \end{beamercolorbox}%
  \par\vspace{0.6em}%
}

\newenvironment{propositionblock}[1]{%
  \par\vspace{0.35em}%
  \begin{beamercolorbox}[
      wd=\textwidth,sep=0.4em,leftskip=0.3em,rightskip=0.6em,
      rounded=false,shadow=false
    ]{block title proposition}%
    \raisebox{0.15em}{\usebeamerfont{block title example}#1}%
  \end{beamercolorbox}%
  \vspace{-0.4em}%
  \begin{beamercolorbox}[
      wd=\textwidth,sep=0.8em,leftskip=0.6em,rightskip=0.6em,
      rounded=false,shadow=false
    ]{block body proposition}%
    \usebeamerfont{block body example}%
}{%
  \end{beamercolorbox}%
  \par\vspace{0.55em}%
}

\newenvironment{illustrationblock}[1]{%
  \par\vspace{0.35em}%
  \begin{beamercolorbox}[
      wd=\textwidth,sep=0.4em,leftskip=0.3em,rightskip=0.6em,
      rounded=false,shadow=false
    ]{block title illustration}%
    \raisebox{0.15em}{\usebeamerfont{block title example}#1}%
  \end{beamercolorbox}%
  \vspace{-0.4em}%
  \begin{beamercolorbox}[
      wd=\textwidth,sep=0.8em,leftskip=0.6em,rightskip=0.6em,
      rounded=false,shadow=false
    ]{block body illustration}%
    \usebeamerfont{block body example}%
}{%
  \end{beamercolorbox}%
  \par\vspace{0.55em}%
}

\newenvironment{proofblock}[1]{%
  \par\vspace{0.3em}%
  \begin{beamercolorbox}[
      wd=\textwidth,sep=0.35em,leftskip=0.3em,rightskip=0.6em,
      rounded=false,shadow=false
    ]{block title proof}%
    \raisebox{0.15em}{\usebeamerfont{block title example}#1}%
  \end{beamercolorbox}%
  \vspace{-0.4em}%
  \begin{beamercolorbox}[
      wd=\textwidth,sep=0.7em,leftskip=0.6em,rightskip=0.6em,
      rounded=false,shadow=false
    ]{block body proof}%
    \usebeamerfont{block body example}%
}{%
  \end{beamercolorbox}%
  \par\vspace{0.45em}%
}


% Self-contained Deck 0. Theme, palette, and teaching-block macros above
% are inherited verbatim from the Fall 2026 FRE6233 Lecture 1 source.
% C++ listings follow the FRE6883 lecture conventions: explicit types,
% std:: qualification, braces on separate lines, and short // comments.
\usepackage{listings}
\lstdefinestyle{coursecpp}{
  language=C++,
  basicstyle=\ttfamily\normalsize,
  keywordstyle=\color{blue},
  commentstyle=\color{MatrixGreen},
  stringstyle=\color{Crimson},
  showstringspaces=false,
  columns=fullflexible,
  keepspaces=true,
  tabsize=4,
  breaklines=false,
  frame=none,
  aboveskip=0.5em,
  belowskip=0.5em
}
\lstset{style=coursecpp}
\setbeamercovered{invisible}

\title{FRE6883 -- Financial Computation}
\subtitle{C++ Fundamentals, Binomial Pricing, and Git}
\author{}
\date{}

% DELIVERY: 8 frames, 12 PDF pages (four answer overlays), about 14 minutes.
% Pacing: 0:30 / 1:30 / 2:00 / 2:00 / 2:00 / 2:00 / 2:00 / 2:00.
% Before each reveal: think silently, vote, invite one explanation, then click.
% Count "not sure" responses as useful diagnostic information, not errors.
% ADAPTATION: if roughly one third struggle with copying or the loop, use
% line-by-line tracing and starter skeletons in Decks 1, 2, and 4. If these
% are comfortable, spend more time on tests and independent repairs.
% The reference item diagnoses prior Topic 2 exposure only; it must not
% become a prerequisite for the rest of today's session. Reveal briefly
% without a reference tutorial if the notation is unfamiliar to most.
% Integer-division uncertainty calls for explicit operand-type checks.
% Keep individual votes before pair discussion so confidence is observable.

\begin{document}

% 1. Welcome (0:30).
\begin{frame}[plain]
  \vfill
  {\usebeamerfont{title}\usebeamercolor[fg]{title}\inserttitle\par}
  \medskip
  {\color{BootcampBlueDark}\rule{0.95\textwidth}{0.8pt}}
  \bigskip
  {\Large\bfseries Recitation 1\par}
  \medskip
  {\large\insertsubtitle\par}
  \vfill
\end{frame}

% 2. Purpose (1:30). Welcome different levels of experience explicitly.
\begin{frame}{What are recitations for?}

\begin{itemize}\setlength{\itemsep}{0.65em}
    \item \textbf{Complement the lectures} with additional explanations, examples, and practice.
    
    \item \textbf{Review the week's material}, focusing on the concepts and techniques introduced in class.
    
    \item \textbf{Ask questions} about anything that was unclear or deserves more practice.
    
    \item \textbf{Practice together} through exercises, coding examples, and debugging.
\end{itemize}

\medskip

\begin{illustrationblock}{Come with your questions!}
Recitations are here to support you throughout the course. 
Please bring any questions you have from the week's lectures, homework, or coding exercises.

\medskip
You are also welcome to \textbf{send me your questions in advance} so that I can prepare examples or exercises around them.
\end{illustrationblock}

\end{frame}
% 3. Background (2:00). Categories overlap: select all that apply.
% Also count Topic 2 exposure separately from general programming experience.
\begin{frame}{What is your programming background?}
  \textbf{Tell me which description fits.}
  \medskip
  \begin{enumerate}\setlength{\itemsep}{0.25em}
    \item Never programmed before
    \item Python / R / Matlab only
    \item Some C / C++
    \item Took a full C++ course
    \item Used C++ in a project or internship
    \item Comfortable with pointers, references, and classes
  \end{enumerate}

\end{frame}

% 4. Copy semantics (2:00). Snippet is inside main with <iostream> included.
\begin{frame}[fragile]{Predict the output: copying a value}
  \relax
  {\small Assume \texttt{<iostream>} is included; this code is inside \texttt{main()}.}
  \medskip
  \begin{columns}[T,onlytextwidth]
    \begin{column}{0.57\textwidth}
\begin{lstlisting}
int x = 3;
int y = x;
y++;
std::cout << x << " " << y;
\end{lstlisting}
    \end{column}
    \begin{column}{0.39\textwidth}
      \begin{overlayarea}{\linewidth}{4.8cm}
        \only<1>{
          \begin{illustrationblock}{Think, then vote}
            \textbf{A}\quad \texttt{3 3}\par\smallskip
            \textbf{B}\quad \texttt{3 4}\par\smallskip
            \textbf{C}\quad \texttt{4 4}\par\smallskip
            \textbf{D}\quad Not sure
          \end{illustrationblock}
          Does changing \texttt{y} change \texttt{x}?
        }
        \only<2>{
          \begin{propositionblock}{Answer B: \texttt{3 4}}
            \texttt{y} receives a copy of the value in \texttt{x}.\par\medskip
            \texttt{y++} increases \texttt{y} by one; \texttt{x} stays 3.
          \end{propositionblock}
        }
      \end{overlayarea}
    \end{column}
  \end{columns}
\end{frame}

% 5. Loop (2:00). Ask for the actual i values before explaining the bound.
\begin{frame}[fragile]{Predict the output: a short loop}
  \relax
  {\small Assume \texttt{<iostream>} is included; this code is inside \texttt{main()}.}
  \medskip
  \begin{columns}[T,onlytextwidth]
    \begin{column}{0.57\textwidth}
\begin{lstlisting}
for (int i = 0; i < 3; i++)
{
    std::cout << i << " ";
}
\end{lstlisting}
    \end{column}
    \begin{column}{0.39\textwidth}
      \begin{overlayarea}{\linewidth}{4.8cm}
        \only<1>{
          \begin{illustrationblock}{Think, then vote}
            \textbf{A}\quad \texttt{0 1 2}\par\smallskip
            \textbf{B}\quad \texttt{0 1 2 3}\par\smallskip
            \textbf{C}\quad \texttt{1 2 3}\par\smallskip
            \textbf{D}\quad Not sure
          \end{illustrationblock}
          How many times does the body run?
        }
        \only<2>{
          \begin{propositionblock}{Answer A: \texttt{0 1 2}}
            The body runs three times.\par\medskip
            At \texttt{i = 3}, \texttt{i < 3} is false, so the loop stops.
          \end{propositionblock}
          {\small A space follows each printed number.}
        }
      \end{overlayarea}
    \end{column}
  \end{columns}
\end{frame}

% 6. Optional reference diagnostic (2:00). Do not infer beginner failure.
% Includes are assumed only to keep the complete function/main pair short.
\begin{frame}[fragile]{What happens? A reference parameter}
  \relax
  {\small\textbf{Optional diagnostic:} new notation is welcome. Assume \texttt{<iostream>} is included.}
  \smallskip
  \begin{columns}[T,onlytextwidth]
    \begin{column}{0.57\textwidth}
\begin{lstlisting}
void f(int& x)
{
    x++;
}
int main()
{
    int a = 3;
    f(a);
    std::cout << a;
    return 0;
}
\end{lstlisting}
    \end{column}
    \begin{column}{0.39\textwidth}
      \begin{overlayarea}{\linewidth}{4.8cm}
        \only<1>{
          \begin{illustrationblock}{Think, then vote}
            \textbf{A}\quad Prints \texttt{3}\par\smallskip
            \textbf{B}\quad Prints \texttt{4}\par\smallskip
            \textbf{C}\quad Does not compile\par\smallskip
            \textbf{D}\quad Not sure
          \end{illustrationblock}
          Is \texttt{x} a copy of \texttt{a} here?
        }
        \only<2>{
          \begin{propositionblock}{Answer B: prints \texttt{4}}
            \texttt{int\& x} makes \texttt{x} a reference to \texttt{a}: another name for the same variable.\par\medskip
            Increasing \texttt{x} changes \texttt{a}.
          \end{propositionblock}
        }
      \end{overlayarea}
    \end{column}
  \end{columns}
\end{frame}

% 7. Integer division (2:00). Keep the operands visible during explanation.
\begin{frame}[fragile]{Predict the output: integer division}
  \relax
  {\small Assume \texttt{<iostream>} is included; this code is inside \texttt{main()}.}
  \medskip
  \begin{columns}[T,onlytextwidth]
    \begin{column}{0.57\textwidth}
\begin{lstlisting}
int x = 5;
double y = x / 2;
std::cout << y;
\end{lstlisting}
    \end{column}
    \begin{column}{0.39\textwidth}
      \begin{overlayarea}{\linewidth}{4.8cm}
        \only<1>{
          \begin{illustrationblock}{Think, then vote}
            \textbf{A}\quad \texttt{2}\par\smallskip
            \textbf{B}\quad \texttt{2.5}\par\smallskip
            \textbf{C}\quad Does not compile\par\smallskip
            \textbf{D}\quad Not sure
          \end{illustrationblock}
          Does declaring \texttt{y} as \texttt{double} change the division?
        }
        \only<2>{
          \begin{propositionblock}{Answer A: prints \texttt{2}}
            Both operands of \texttt{x / 2} are integers: the result is 2.\par\smallskip
            It is then stored as a \texttt{double}.\par\smallskip
            Use \texttt{x / 2.0} to obtain 2.5.
          \end{propositionblock}
        }
      \end{overlayarea}
    \end{column}
  \end{columns}
\end{frame}

% 8. Transition (2:00). Invite a plain-language answer, not C++ syntax yet.
% Move into Deck 1 with the scaffolding indicated by the votes above.
\begin{frame}{From a financial problem to an algorithm}
  \begin{definitionblock}{A European call at expiry}
    Given the terminal stock price $z$ and strike $K$,\par
    compute the payoff: $z-K$ if $z>K$, and zero otherwise.
  \end{definitionblock}
  \medskip
  \textbf{Before writing C++:}
  \begin{enumerate}\setlength{\itemsep}{0.35em}
    \item What are the inputs and the output?
    \item What steps should the program follow?
    \item Which examples would convince you it works?
  \end{enumerate}
  \medskip
  {\color{BootcampBlueDark}\textbf{Next: describe the steps, then translate them into code.}}
\end{frame}

\end{document}
